%%
%% Speech Analytics e Processamento de Linguagem Natural Aplicados à
%% Comunicação Clínica: Uma Revisão Bibliográfica
%%
%% Autores: Rafael Baena Neto, Vicente Idalberto Becerra Sablón,
%%          Fabrício Evangelista dos Santos, Lucas Kenzo Terazzin Ida,
%%          Murilo Nogueira Minutti
%%
%% Periódico-alvo: International Journal of Medical Informatics (IJMI)
%% APC coberto: Acordo CAPES/Elsevier Read & Publish 2026–2028
%% Classe: cas-sc.cls v2.4 (coluna única — para submissão/revisão)
%%
%% NOTA: Este arquivo está em português para revisão interna.
%% Tradução para o inglês será realizada na etapa final (Gate 4).
%%

\documentclass[a4paper,fleqn]{cas-sc}

\usepackage[authoryear,longnamesfirst]{natbib}
\usepackage[utf8]{inputenc}
\usepackage[T1]{fontenc}
\usepackage{lmodern}
\usepackage[brazil]{babel}
\usepackage{amsmath}
\usepackage{amssymb}
\DeclareSymbolFont{CMSYfix}{OMS}{cmsy}{m}{n}
\DeclareMathSymbol{\approx}{\mathrel}{CMSYfix}{"19}
\usepackage{csquotes}
\usepackage{graphicx}
\usepackage{booktabs}
\usepackage{longtable}
\usepackage{array}
\usepackage{hyperref}
\usepackage{xcolor}
\usepackage{tikz}

\hypersetup{colorlinks=true, linkcolor=blue, citecolor=blue, urlcolor=blue}

%% ================================================================
\begin{document}
\let\WriteBookmarks\relax
\def\floatpagepagefraction{1}
\def\textpagefraction{.001}

%% ----------------------------------------------------------------
%% Cabeçalho corrente
%% ----------------------------------------------------------------
\shorttitle{Speech Analytics na Comunicação Clínica}
\shortauthors{R. Baena Neto et~al.}

%% ----------------------------------------------------------------
%% Título
%% ----------------------------------------------------------------
\title[mode=title]{Speech Analytics e Processamento de Linguagem Natural
Aplicados à Comunicação Clínica: Uma Revisão Bibliográfica}

%% ----------------------------------------------------------------
%% Autores e afiliações
%% Afiliação [1]: Pós-Graduação em Ciência de Dados em Saúde, USF
%% Afiliação [2]: Engenharia da Computação, USF (alunos de IC)
%% ----------------------------------------------------------------
\author[1]{Rafael Baena Neto}[%
  type=editor,
  orcid=0009-0005-2122-4665,
  role=Researcher
]
\cormark[1]
\fnmark[1]
\ead{rafael.baena@mail.usf.edu.br}
\credit{Conceptualização, Metodologia, Curadoria de dados, Análise formal,
        Redação -- rascunho original}

\affiliation[1]{%
  organization={Programa de Pós-Graduação em Ciência de Dados em Saúde,
                Universidade São Francisco (USF)},
  addressline={Rua Alexandre Rodrigues Barbosa, 45},
  city={Bragança Paulista},
  postcode={12916-900},
  state={SP},
  country={Brasil}
}

\author[1]{Vicente Idalberto Becerra Sablón}[%
  orcid=0000-0003-3127-1906
]
\fnmark[2]
\ead{vicente.sablon@usf.edu.br}
\credit{Supervisão, Conceptualização, Redação -- revisão e edição}

\author[2]{Fabrício Evangelista dos Santos}[%
  orcid=0009-0005-8163-3479
]
\fnmark[3]
\ead{fabricio.evangelista@mail.usf.edu.br}
\credit{Investigação, Curadoria de dados, Metodologia, Software,
        Redação -- revisão e edição}

\author[2]{Lucas Kenzo Terazzin Ida}[%
  orcid=0009-0009-2328-5022
]
\fnmark[4]
\ead{lucas.ida@mail.usf.edu.br}
\credit{Investigação, Curadoria de dados, Análise formal,
        Redação -- revisão e edição}

\author[2]{Murilo Nogueira Minutti}[%
  orcid=0009-0000-6411-1576
]
\fnmark[5]
\ead{murilo.minutti@mail.usf.edu.br}
\credit{Investigação, Curadoria de dados, Análise formal, Visualização,
        Redação -- revisão e edição}

\affiliation[2]{%
  organization={Curso de Engenharia da Computação,
                Universidade São Francisco (USF)},
  addressline={Rua Alexandre Rodrigues Barbosa, 45},
  city={Bragança Paulista},
  postcode={12916-900},
  state={SP},
  country={Brasil}
}

\cortext[1]{Autor correspondente}
\fntext[1]{Mestrando, Programa de Pós-Graduação em Ciência de Dados em Saúde, USF.}
\fntext[2]{Docente e orientador, Programa de Pós-Graduação em Ciência de Dados em Saúde, USF.}
\fntext[3]{Aluno de Iniciação Científica, Engenharia da Computação, USF.}
\fntext[4]{Aluno de Iniciação Científica, Engenharia da Computação, USF.}
\fntext[5]{Aluno de Iniciação Científica, Engenharia da Computação, USF.}

%% ----------------------------------------------------------------
%% Resumo — máximo 250 palavras (IJMI)
%% ----------------------------------------------------------------
\begin{abstract}
\textbf{Contexto:}
A comunicação profissional-paciente constitui uma fonte densa de informação
clínica que ainda permanece predominantemente não estruturada e subutilizada
analiticamente. Avanços recentes em Reconhecimento Automático de Fala (ASR),
Processamento de Linguagem Natural (NLP) e análise de biomarcadores vocais
têm ampliado a capacidade de extrair dados estruturados a partir de interações
clínicas na anamnese. Ainda assim, poucos estudos operacionalizam essas
interações como variáveis mensuráveis derivadas da fala, especialmente no que
diz respeito a marcadores comunicacionais e comportamentais que podem ser
utilizados para análise de dados e avaliação de resultados no cuidado do
paciente.

\textbf{Objetivo:}
Identificar, sintetizar e avaliar criticamente as evidências sobre a aplicação
de \textit{Speech Analytics} e NLP à comunicação clínica profissional-paciente,
mapeando abordagens metodológicas, marcadores identificados, desafios técnicos
e lacunas para pesquisa futura.

\textbf{Métodos:}
Revisão bibliográfica seguindo as diretrizes PRISMA-ScR \citep{Tricco2018}
e JBI \citep{Peters2020}. A busca foi expandida para múltiplas bases de dados eletrônicas indexadas (PubMed/MEDLINE, Scopus, Web of Science, CINAHL/EBSCO, Cochrane Library e Google Scholar) sem filtros de texto completo gratuito, e complementada por busca manual via \textit{snowballing}. O corpus final compreende 24~artigos selecionados a partir de 5.621~registros da busca eletrônica primária (após remoção de 2.432 duplicatas e triagem em duas etapas) e 7 estudos adicionados manualmente por busca ativa, organizados em quatro eixos temáticos: (A)~NLP clínico e extração de entidades
biomédicas; (B)~documentação clínica automatizada e ASR;
(C)~comunicação médico-paciente; (D)~biomarcadores vocais e análise multimodal.

\textbf{Resultados:}
Os estudos evidenciam avanços consolidados em NLP clínico (modelos baseados
em \textit{transformers}, BERT lidera o Reconhecimento de Entidades Nomeadas)
e em transcrição automatizada (WER $\approx$9\% alcançável com
Whisper/PyAnnote). A comunicação centrada no paciente demonstra impacto
mensurável em desfechos objetivos de saúde. Marcadores comunicacionais autorreparo,
engajamento, suporte emocional e indicadores de barreira são identificados
como preditores de adesão. A análise multimodal
(acústica\,+\,semântica) permanece em estado inicial; a maioria das
abordagens foca em automação documental, negligenciando a interação clínica
como objeto de análise em si.

\textbf{Conclusões:}
Confirma-se uma lacuna consistente: poucos estudos operacionalizam interações
clínicas como variáveis mensuráveis derivadas da fala, especialmente
marcadores comunicacionais e comportamentais extraídos da anamnese. Esta
revisão fundamenta a necessidade de estudos que integrem \textit{Speech
Analytics} à análise estruturada da comunicação clínica como fonte de dados
para monitoramento de adesão e qualidade do cuidado.

\end{abstract}

%% ----------------------------------------------------------------
%% Research Highlights — OBRIGATÓRIO no IJMI (3 itens, máx. 85 caracteres)
%% ----------------------------------------------------------------
\begin{highlights}
\item Modelos BERT lideram NER clínico; WER $\approx$9\% alcançável com
      pipelines Whisper/PyAnnote
\item Marcadores comunicacionais (autorreparo, engajamento, barreiras)
      predizem adesão terapêutica
\item Lacuna confirmada: interações clínicas raramente operacionalizadas
      como variáveis derivadas da fala
\end{highlights}

%% ----------------------------------------------------------------
%% Palavras-chave
%% ----------------------------------------------------------------
\begin{keywords}
Speech analytics \sep
Processamento de linguagem natural \sep
Comunicação clínica \sep
Reconhecimento automático de fala \sep
Marcadores comunicacionais \sep
Digital scribe
\end{keywords}

\maketitle

%% ================================================================
\section{Introdução}
\label{sec:intro}

A comunicação clínica constitui uma das mais ricas e menos exploradas fontes
de dados em saúde. Em sistemas públicos de saúde, anamneses e orientações
terapêuticas, interações verbais com estrutura propedêutica definida
pela semiologia clínica \citep{Porto2019}, ocorrem predominantemente
como trocas não estruturadas do ponto de vista analítico, raramente
convertidas em dados mensuráveis \citep{Wang2018,Koleck2019}.

A qualidade da interação profissional-paciente demonstra impacto mensurável
sobre a adesão e os desfechos clínicos objetivos \citep{Kelley2014}. Barreiras
comunicacionais, incluindo lacunas de literacia em saúde, expectativas
desalinhadas e padrões de miscomunicação, associam-se à baixa adesão
terapêutica e pior prognóstico \citep{McCabe2018,Alkhamees2023}.

O \textit{Speech Analytics}, análise computacional da linguagem falada,
oferece um caminho para extrair indicadores estruturados e mensuráveis
dessas interações. Avanços recentes em ASR, especialmente modelos baseados
em \textit{transformers} como o Whisper \citep{Radford2023}, combinados a
grandes modelos de linguagem (LLMs) com instruções, reduziram as barreiras
para implantação de \textit{pipelines} de NLP ponta-a-ponta na saúde
\citep{Klusty2025}. Ainda assim, a maioria das abordagens existentes visa
suporte diagnóstico ou automação documental, não tratando a interação clínica
em si como objeto de análise estruturada \citep{vanBuchem2021}. O desenvolvimento
de um protocolo de \textit{Speech Analytics} para monitoramento da reabilitação
ortopédica em sistema público de saúde \citep{BaenaNeto2026} exemplifica o
potencial translacional dessas tecnologias e evidencia a necessidade de um
mapa consolidado de evidências para subsidiar o desenho de estudos prospectivos.

Revisões recentes têm abordado aspectos parcialmente sobrepostos a este
escopo. \citet{Falcetta2023} identificaram ausência
de validação prospectiva em sistemas de documentação automática de
interações clínicas. \citet{Ng2025} avaliaram o desempenho
técnico de sistemas de ASR para documentação clínica, reportando WER
entre 8,7\% e mais de 50\% conforme o contexto. \citet{Alabbad2026}
mapearam aplicações de ASR na saúde na era
pós-LLM. \citet{Albert2024} realizaram uma revisão de
escopo sobre métodos de IA, fala e NLP para avaliação da comunicação
clínico-paciente. Nenhuma dessas revisões, contudo, concentra-se na
operacionalização da interação clínica como conjunto de \textbf{variáveis
mensuráveis derivadas da fala}, especialmente marcadores comunicacionais
e comportamentais (autorreparo, engajamento, barreiras, literacia), como
objeto de análise em si, independente da documentação ou do diagnóstico.
É precisamente essa lacuna que a presente revisão se propõe a mapear.

Esta revisão bibliográfica tem como objetivo mapear as evidências sobre
\textit{Speech Analytics} e NLP aplicados à comunicação clínica,
caracterizando o panorama metodológico e identificando lacunas para
pesquisas futuras.

%% ================================================================
\section{Métodos}
\label{sec:metodos}

\subsection{Desenho do estudo}

Esta revisão bibliográfica segue o referencial metodológico proposto por
\citet{Arksey2005}, refinado por \citet{Levac2010} e pelo Instituto Joanna
Briggs (JBI) \citep{Peters2020}. A condução e o relato seguem as diretrizes
PRISMA-ScR (\textit{Preferred Reporting Items for Systematic Reviews and
Meta-Analyses} extensão para Scoping Reviews) \citep{Tricco2018}.

\subsection{Perguntas de pesquisa}

\begin{itemize}
  \item \textbf{Q1:} Quais abordagens de \textit{Speech Analytics} têm sido
        aplicadas à comunicação clínica nos últimos cinco anos?
  \item \textbf{Q2:} Quais métodos de NLP e ASR são utilizados para transcrição
        e análise automatizada de consultas clínicas?
  \item \textbf{Q3:} Quais marcadores comunicacionais e comportamentais foram
        identificados ou propostos na interação profissional-paciente?
  \item \textbf{Q4:} Quais são os desafios técnicos, éticos e metodológicos
        para a análise automatizada da comunicação clínica?
\end{itemize}

\subsection{Critérios de elegibilidade}

\subsubsection{Critérios de inclusão}
\begin{itemize}
  \item \textbf{IC1:} O estudo aplica tecnologias de \textit{Speech Analytics},
        NLP ou ASR para transcrição clínica e extração de dados médicos não
        estruturados, ou analisa e mensura os fundamentos da comunicação
        profissional-paciente.
  \item \textbf{IC2:} O estudo foca em interações baseadas em texto livre de
        prontuários (\textit{free-text}), diálogos clínicos falados (consultas),
        geração de documentação automatizada (escribas digitais/notas SOAP) ou
        marcos teóricos da tomada de decisão compartilhada.
  \item \textbf{IC3:} Disponibilidade de texto completo em inglês, português ou
        espanhol; publicação no período de janeiro de 2014 a 2026.
  \item \textbf{IC4:} Apresenta desenvolvimento original, teste de viabilidade
        empírica, validação de algoritmos, ou constitui revisão sistemática ou de
        escopo anterior sobre o tema.
\end{itemize}

\subsubsection{Critérios de exclusão}
\begin{itemize}
  \item \textbf{EC1:} Estudos que não abordem PLN, ASR ou análise de áudio de
        interações clínicas.
  \item \textbf{EC2:} Pesquisas que descrevam sistemas de transcrição ou extração
        sem analisar a qualidade da interação ou o impacto na relação
        profissional-paciente.
  \item \textbf{EC3:} Foco exclusivo em diagnóstico por voz de condições físicas
        ou psiquiátricas sem análise da interação dialógica.
  \item \textbf{EC4:} Artigos de opinião, comentários, editoriais, resumos de
        conferência ou livros.
  \item \textbf{EC5:} Estudos que não descrevem coleta de dados, metodologia de
        processamento ou validação técnica.
  \item \textbf{EC6:} Pesquisas que abordem exclusivamente interações não clínicas
        ou contextos simulados sem transferibilidade prática demonstrada.
\end{itemize}

\subsection{Fontes de informação e estratégia de busca}

A busca bibliográfica foi conduzida em múltiplas bases de dados eletrônicas indexadas — PubMed/MEDLINE, Scopus, Web of Science, CINAHL/EBSCO, Cochrane Library e Google Scholar —, visando cobrir de forma abrangente as interseções entre saúde, processamento de linguagem natural e análise de fala, segundo as diretrizes do PRISMA-ScR \citep{Tricco2018} e do Joanna Briggs Institute \citep{Peters2020}. 

A base de dados primária foi o PubMed/MEDLINE, na qual foram aplicadas todas as seis estratégias de busca formuladas (cinco strings especializadas nos eixos temáticos e uma string geral unificada). Para as bases Scopus, Web of Science e CINAHL/EBSCO, executaram-se as cinco strings especializadas. Na Cochrane Library, a busca concentrou-se na String 1 adaptada para campos de título, resumo e palavras-chave (\texttt{:ti,ab,kw}), adequada ao foco da base em ensaios clínicos e intervenções. No Google Scholar, aplicou-se o operador \texttt{allintitle:} nas cinco strings especializadas para respeitar a restrição de 256 caracteres do mecanismo.

Foram estabelecidos como filtros eletrônicos gerais o período de publicação entre janeiro de 2014 e 2026, e artigos nos idiomas inglês, espanhol e português. O filtro de texto completo gratuito (\textit{free full text}) foi removido para evitar viés de seleção.

Cada pesquisador utilizou mais de uma string de busca para realizar as pesquisas. As strings específicas utilizadas foram as seguintes:

\begin{enumerate}
  \item \textbf{String 1 (PLN, Sintomas e Prontuários):}
        \begin{verbatim}
("natural language processing"[Title/Abstract])
AND
(symptoms[Title/Abstract] OR "patient-reported symptoms"[Title/Abstract])
AND
("electronic health records"[Title/Abstract] OR "patient-centered"[Title/Abstract])
        \end{verbatim}

  \item \textbf{String 2 (Extração de Informação Clínica e NER):}
        \begin{verbatim}
("clinical information extraction"[Title/Abstract]
OR "medical information extraction"[Title/Abstract]
OR "named entity recognition"[Title/Abstract])
AND
("medical free text"[Title/Abstract]
OR "review"[Title/Abstract]
OR "patient-generated health data"[Title/Abstract])
        \end{verbatim}

  \item \textbf{String 3 (Comunicação Clínico-Paciente e Desfechos):}
        \begin{verbatim}
("patient-physician communication"[Title/Abstract]
OR "doctor-patient communication"[Title/Abstract]
OR "patient-clinician relationship"[Title/Abstract])
AND
("intercultural settings"[Title/Abstract]
OR miscommunication[Title/Abstract]
OR "healthcare outcomes"[Title/Abstract]
OR "shared decision-making"[Title/Abstract]
OR "communication quality scale"[Title/Abstract])
        \end{verbatim}

  \item \textbf{String 4 (Biomarcadores Vocais e Speech Analytics):}
        \begin{verbatim}
("vocal biomarkers"[Title/Abstract] OR "speech analytics"[Title/Abstract])
AND
("clinical practice"[Title/Abstract]
OR "voice for health"[Title/Abstract]
OR "SUS"[Title/Abstract])
        \end{verbatim}

  \item \textbf{String 5 (NLP, Speech e Revisões):}
        \begin{verbatim}
(("Applying natural language processing"[Title/Abstract])
  AND ("Systematic review"[Title/Abstract]))
OR
(("Identifying Topics"[Title/Abstract])
  AND ("Natural Language Processing Methods"[Title/Abstract]))
OR
(("Automatic speech recognition"[Title/Abstract])
  AND ("patient-clinician conversations"[Title/Abstract]))
        \end{verbatim}
\end{enumerate}

Adicionalmente, foi executada uma busca abrangente utilizando a seguinte string geral unificada:

\begin{verbatim}
("speech analytics" OR "speech analysis" OR "speech recognition"
OR "automatic speech recognition" OR ASR OR "digital scribe"
OR "clinical transcription" OR "clinical documentation" OR "vocal biomarkers"
OR "healthcare conversations" OR "doctor-patient conversations"
OR "patient-doctor conversations")
AND
("natural language processing" OR NLP OR "large language model" OR LLM
OR "information extraction" OR "clinical information extraction"
OR "named entity recognition" OR NER OR "relation extraction" OR "text mining"
OR "conversation analysis" OR "topic modeling" OR "SOAP note" OR "SOAP notes")
AND
(healthcare OR clinical OR medicine OR medical OR physician OR clinician
OR doctor OR patient OR "electronic health record" OR EHR OR hospital
OR outpatient)
\end{verbatim}

Para mitigar a perda de artigos relevantes e enriquecer qualitativamente o corpus, a busca eletrônica foi complementada pelo método de \textit{snowballing} (rastreamento manual de referências ou \textit{backward snowballing}). A partir das revisões sistemáticas e dos artigos fundamentais localizados na busca primária, realizou-se a análise manual de suas listas de referências para identificar e incluir outros estudos que atendem estritamente aos critérios de elegibilidade definidos, assegurando a cobertura de trabalhos chaves de relevância metodológica ou conceitual para a presente revisão.

\subsection{Seleção de estudos}

Os registros recuperados da busca primária no PubMed/MEDLINE foram importados no Rayyan
\citep{Ouzzani2016} para deduplicação e triagem. Os títulos e resumos foram
triados de forma independente por dois revisores (F.E.S. e M.M.), com
concordância avaliada pelo coeficiente kappa de Cohen ($\kappa = 0{,}78$). Desacordos
foram resolvidos por um terceiro revisor (L.K.T.I.) ou pelo autor correspondente.
Os textos completos dos registros potencialmente elegíveis foram então
recuperados e avaliados conforme os critérios de elegibilidade. O processo de
seleção está documentado no fluxograma PRISMA-ScR (Fig.~\ref{fig:prisma}).
%% κ = 0,78 (F.E.S. e M.M.; triagem dupla-cega Rayyan → irr::kappa2, weight="unweighted")

\subsection{Extração de dados}

Os dados foram extraídos com formulário padronizado contendo: autor(es), ano,
país, desenho do estudo, contexto clínico, características da população,
tecnologia utilizada (modelo ASR, arcabouço NLP, LLM), variáveis clínicas
extraídas, desfechos reportados, arcabouço ético e limitações.

\subsection{Síntese}

Foi realizada uma síntese narrativa. Os achados estão organizados em torno
dos quatro eixos temáticos definidos nos critérios de elegibilidade:
(A)~NLP clínico e extração de entidades biomédicas;
(B)~documentação clínica automatizada e ASR;
(C)~comunicação médico-paciente; e
(D)~biomarcadores vocais e análise multimodal.

%% ================================================================
\section{Resultados}
\label{sec:resultados}

% ---------------------------------------------------------------
\subsection{Seleção de estudos}
% ---------------------------------------------------------------

As buscas nas bases de dados eletrônicas (PubMed/MEDLINE, Scopus, Web of Science, CINAHL/EBSCO, Cochrane Library e Google Scholar) recuperaram 5.621 registros. Após o processo de deduplicação realizado no Zotero, que removeu 2.432 registros duplicados, obteve-se um total de 3.189 registros únicos. Estes foram submetidos à triagem por título, resultando na exclusão de 2.680 artigos pelos critérios de elegibilidade (1.120 por EC1, 584 por EC2, 134 por EC3, 281 por EC4, 358 por EC5 e 203 por EC6) e na seleção de 509 artigos aprovados para a triagem por resumo. Na triagem por resumo, foram excluídos 474 artigos (183 por EC1, 116 por EC2, 29 por EC3, 43 por EC4, 61 por EC5 e 42 por EC6), resultando em 35 artigos selecionados para avaliação de texto completo. Adicionalmente, 7 artigos identificados via \textit{snowballing} (rastreamento manual de referências) foram incluídos na avaliação de elegibilidade. Na fase de leitura na íntegra de texto completo (35 das bases e 7 do snowballing), foram excluídos 18 artigos das bases pelos critérios de elegibilidade (4 por EC1, 3 por EC2, 3 por EC3, 2 por EC4, 4 por EC5 e 2 por EC6) e 0 artigos de snowballing, restando 17 artigos das bases e 7 de snowballing. Ao final, 24 estudos foram incluídos na síntese qualitativa (17 da busca eletrônica nas bases e 7 provenientes do rastreamento manual). O fluxograma PRISMA-ScR detalhando todo o processo está apresentado na Fig.~\ref{fig:prisma}.

\begin{figure}[h]
  \centering
  \resizebox{\linewidth}{!}{%% prisma_flow.tex — PRISMA-ScR Flowchart
%% Included via %% prisma_flow.tex — PRISMA-ScR Flowchart
%% Included via %% prisma_flow.tex — PRISMA-ScR Flowchart
%% Included via \input{prisma_flow} in the figure environment of main.tex
%% Exact counts: 5,621 identified across 5 databases, 2,432 duplicates removed,
%% 3,189 screened by title (2,680 excluded), 509 screened by abstract
%% (474 excluded), 35 full text (databases) + 7 snowballing, 18 excluded,
%% 24 included (17 databases + 7 snowballing).

\usetikzlibrary{positioning, calc, arrows.meta}

\begin{tikzpicture}[
  %% main flow boxes
  mbox/.style={
    rectangle, draw=black, thick, rounded corners=2pt,
    minimum width=6.8cm, text width=6.6cm,
    minimum height=1.3cm, align=center,
    fill=white, font=\small
  },
  %% exclusion boxes (right margin)
  ebox/.style={
    rectangle, draw=black, thick, rounded corners=2pt,
    minimum width=5.0cm, text width=4.8cm,
    minimum height=1.3cm, align=flush left,
    fill=gray!8, font=\small
  },
  %% arrows
  arr/.style={-{Stealth[length=5pt,width=4pt]}, thick},
  %% phase labels
  phase/.style={font=\bfseries\scriptsize, rotate=90, anchor=center}
]

%% ================================================================
%% PHASE 1 — IDENTIFICATION
%% ================================================================
\node[mbox, minimum width=6.0cm, text width=5.8cm] (dbbox) at (-1.8, 0) {
  \textbf{Electronic databases}\\[2pt]
  $n = 5{,}621$\\[2pt]
  {\scriptsize Web of Science: 1,719\\
    Scopus: 1,580\\
    PubMed/MEDLINE: 1,059\\
    EBSCO/CINAHL: 586\\
    Cochrane Library: 399\\
    Google Scholar: 278}
};

\node[mbox, minimum width=5.0cm, text width=4.8cm,
      right=1.5cm of dbbox] (otherbox) {
  \textbf{Other sources (Snowballing)}\\[2pt]
  $n = 7$\\[2pt]
  {\scriptsize (Manual backward reference\\
    tracking)}
};

%% ================================================================
%% PHASE 2 — SCREENING
%% ================================================================
\coordinate (midid) at ($(dbbox.south)!0.5!(otherbox.south)$);

\node[mbox, below=2.6cm of midid, anchor=north] (dedup) {
  \textbf{Database records after deduplication}\\[2pt]
  $n = 3{,}189$\quad{\scriptsize (5,621 $-$ 2,432 duplicates)}
};

%% arrow from databases to deduplication
\draw[arr] (dbbox.south) -- ++(0,-0.8cm) -| (dedup.north);

\node[mbox, below=1.2cm of dedup] (screen1) {
  \textbf{Title screening}\\[2pt]
  $n = 3{,}189$
};

\draw[arr] (dedup) -- (screen1);

\node[ebox, right=2.2cm of screen1] (excl1) {
  \textbf{Excluded by title:} $n = 2{,}680$\\[3pt]
  {\scriptsize
    EC1 (No NLP/ASR of clinical interactions): 1,120\\
    EC2 (No interaction-quality analysis): 584\\
    EC3 (Voice-based diagnosis only): 134\\
    EC4 (Publication type): 281\\
    EC5 (No methodological description): 358\\
    EC6 (Simulated/non-clinical context): 203
  }
};

\draw[arr] (screen1.east) -- (excl1.west);

\node[mbox, below=1.2cm of screen1] (screen2) {
  \textbf{Abstract screening}\\[2pt]
  $n = 509$
};

\draw[arr] (screen1) -- (screen2);

\node[ebox, right=2.2cm of screen2] (excl2) {
  \textbf{Excluded by abstract:} $n = 474$\\[3pt]
  {\scriptsize
    EC1 (No NLP/ASR of clinical interactions): 183\\
    EC2 (No interaction-quality analysis): 116\\
    EC3 (Voice-based diagnosis only): 29\\
    EC4 (Publication type): 43\\
    EC5 (No methodological description): 61\\
    EC6 (Simulated/non-clinical context): 42
  }
};

\draw[arr] (screen2.east) -- (excl2.west);

%% ================================================================
%% PHASE 3 — ELIGIBILITY
%% ================================================================
\node[mbox, below=1.2cm of screen2] (eligib) {
  \textbf{Full text assessed}\\[2pt]
  $n = 35$ (databases) $+$ $7$ (snowballing)
};

\draw[arr] (screen2) -- (eligib);

%% snowballing enters directly at eligibility (bypasses screening)
\draw[arr] (otherbox.south) -- ++(0,-0.8cm) -| ([xshift=1.1cm]eligib.east) -- (eligib.east);

\node[ebox, right=2.2cm of eligib] (excl3) {
  \textbf{Excluded in full text:} $n = 18$ (databases); $0$ (snowballing)\\[3pt]
  {\scriptsize
    Assessed against the full eligibility criteria (EC1--EC6)
  }
};

\draw[arr] (eligib.east) -- (excl3.west);

%% ================================================================
%% PHASE 4 — INCLUDED
%% ================================================================
\node[mbox, below=1.2cm of eligib] (included) {
  \textbf{Studies included in synthesis}\\[2pt]
  $n = 24$\\[2pt]
  {\scriptsize (17 databases + 7 snowballing)}
};

\draw[arr] (eligib) -- (included);

%% ================================================================
%% PHASE LABELS (left margin)
%% ================================================================
\node[phase] at ($(dbbox.west) + (-1.1cm, 0)$)      {IDENTIFICATION};
\node[phase] at ($(dedup.west)  + (-1.1cm, 0)$)      {SCREENING};
\node[phase] at ($(screen1.west) + (-1.1cm, 0)$)      {};   %% screening continues
\node[phase] at ($(screen2.west) + (-1.1cm, 0)$)      {};   %% screening continues
\node[phase] at ($(eligib.west) + (-1.1cm, 0)$)      {ELIGIBILITY};
\node[phase] at ($(included.west) + (-1.1cm, 0)$)    {INCLUDED};

%% ================================================================
%% FOOTNOTE
%% ================================================================
\node[font=\scriptsize, anchor=north west, text width=12cm]
  at ($(included.south west) + (-1.5cm, -0.5cm)$) {%
  Exact counts from the expanded multi-database search (2026).
  Snowballing assessed with the same eligibility criteria.
  $\kappa = 0.78$ (inter-reviewer agreement; Cohen's kappa).
};

\end{tikzpicture}
 in the figure environment of main.tex
%% Exact counts: 5,621 identified across 5 databases, 2,432 duplicates removed,
%% 3,189 screened by title (2,680 excluded), 509 screened by abstract
%% (474 excluded), 35 full text (databases) + 7 snowballing, 18 excluded,
%% 24 included (17 databases + 7 snowballing).

\usetikzlibrary{positioning, calc, arrows.meta}

\begin{tikzpicture}[
  %% main flow boxes
  mbox/.style={
    rectangle, draw=black, thick, rounded corners=2pt,
    minimum width=6.8cm, text width=6.6cm,
    minimum height=1.3cm, align=center,
    fill=white, font=\small
  },
  %% exclusion boxes (right margin)
  ebox/.style={
    rectangle, draw=black, thick, rounded corners=2pt,
    minimum width=5.0cm, text width=4.8cm,
    minimum height=1.3cm, align=flush left,
    fill=gray!8, font=\small
  },
  %% arrows
  arr/.style={-{Stealth[length=5pt,width=4pt]}, thick},
  %% phase labels
  phase/.style={font=\bfseries\scriptsize, rotate=90, anchor=center}
]

%% ================================================================
%% PHASE 1 — IDENTIFICATION
%% ================================================================
\node[mbox, minimum width=6.0cm, text width=5.8cm] (dbbox) at (-1.8, 0) {
  \textbf{Electronic databases}\\[2pt]
  $n = 5{,}621$\\[2pt]
  {\scriptsize Web of Science: 1,719\\
    Scopus: 1,580\\
    PubMed/MEDLINE: 1,059\\
    EBSCO/CINAHL: 586\\
    Cochrane Library: 399\\
    Google Scholar: 278}
};

\node[mbox, minimum width=5.0cm, text width=4.8cm,
      right=1.5cm of dbbox] (otherbox) {
  \textbf{Other sources (Snowballing)}\\[2pt]
  $n = 7$\\[2pt]
  {\scriptsize (Manual backward reference\\
    tracking)}
};

%% ================================================================
%% PHASE 2 — SCREENING
%% ================================================================
\coordinate (midid) at ($(dbbox.south)!0.5!(otherbox.south)$);

\node[mbox, below=2.6cm of midid, anchor=north] (dedup) {
  \textbf{Database records after deduplication}\\[2pt]
  $n = 3{,}189$\quad{\scriptsize (5,621 $-$ 2,432 duplicates)}
};

%% arrow from databases to deduplication
\draw[arr] (dbbox.south) -- ++(0,-0.8cm) -| (dedup.north);

\node[mbox, below=1.2cm of dedup] (screen1) {
  \textbf{Title screening}\\[2pt]
  $n = 3{,}189$
};

\draw[arr] (dedup) -- (screen1);

\node[ebox, right=2.2cm of screen1] (excl1) {
  \textbf{Excluded by title:} $n = 2{,}680$\\[3pt]
  {\scriptsize
    EC1 (No NLP/ASR of clinical interactions): 1,120\\
    EC2 (No interaction-quality analysis): 584\\
    EC3 (Voice-based diagnosis only): 134\\
    EC4 (Publication type): 281\\
    EC5 (No methodological description): 358\\
    EC6 (Simulated/non-clinical context): 203
  }
};

\draw[arr] (screen1.east) -- (excl1.west);

\node[mbox, below=1.2cm of screen1] (screen2) {
  \textbf{Abstract screening}\\[2pt]
  $n = 509$
};

\draw[arr] (screen1) -- (screen2);

\node[ebox, right=2.2cm of screen2] (excl2) {
  \textbf{Excluded by abstract:} $n = 474$\\[3pt]
  {\scriptsize
    EC1 (No NLP/ASR of clinical interactions): 183\\
    EC2 (No interaction-quality analysis): 116\\
    EC3 (Voice-based diagnosis only): 29\\
    EC4 (Publication type): 43\\
    EC5 (No methodological description): 61\\
    EC6 (Simulated/non-clinical context): 42
  }
};

\draw[arr] (screen2.east) -- (excl2.west);

%% ================================================================
%% PHASE 3 — ELIGIBILITY
%% ================================================================
\node[mbox, below=1.2cm of screen2] (eligib) {
  \textbf{Full text assessed}\\[2pt]
  $n = 35$ (databases) $+$ $7$ (snowballing)
};

\draw[arr] (screen2) -- (eligib);

%% snowballing enters directly at eligibility (bypasses screening)
\draw[arr] (otherbox.south) -- ++(0,-0.8cm) -| ([xshift=1.1cm]eligib.east) -- (eligib.east);

\node[ebox, right=2.2cm of eligib] (excl3) {
  \textbf{Excluded in full text:} $n = 18$ (databases); $0$ (snowballing)\\[3pt]
  {\scriptsize
    Assessed against the full eligibility criteria (EC1--EC6)
  }
};

\draw[arr] (eligib.east) -- (excl3.west);

%% ================================================================
%% PHASE 4 — INCLUDED
%% ================================================================
\node[mbox, below=1.2cm of eligib] (included) {
  \textbf{Studies included in synthesis}\\[2pt]
  $n = 24$\\[2pt]
  {\scriptsize (17 databases + 7 snowballing)}
};

\draw[arr] (eligib) -- (included);

%% ================================================================
%% PHASE LABELS (left margin)
%% ================================================================
\node[phase] at ($(dbbox.west) + (-1.1cm, 0)$)      {IDENTIFICATION};
\node[phase] at ($(dedup.west)  + (-1.1cm, 0)$)      {SCREENING};
\node[phase] at ($(screen1.west) + (-1.1cm, 0)$)      {};   %% screening continues
\node[phase] at ($(screen2.west) + (-1.1cm, 0)$)      {};   %% screening continues
\node[phase] at ($(eligib.west) + (-1.1cm, 0)$)      {ELIGIBILITY};
\node[phase] at ($(included.west) + (-1.1cm, 0)$)    {INCLUDED};

%% ================================================================
%% FOOTNOTE
%% ================================================================
\node[font=\scriptsize, anchor=north west, text width=12cm]
  at ($(included.south west) + (-1.5cm, -0.5cm)$) {%
  Exact counts from the expanded multi-database search (2026).
  Snowballing assessed with the same eligibility criteria.
  $\kappa = 0.78$ (inter-reviewer agreement; Cohen's kappa).
};

\end{tikzpicture}
 in the figure environment of main.tex
%% Exact counts: 5,621 identified across 5 databases, 2,432 duplicates removed,
%% 3,189 screened by title (2,680 excluded), 509 screened by abstract
%% (474 excluded), 35 full text (databases) + 7 snowballing, 18 excluded,
%% 24 included (17 databases + 7 snowballing).

\usetikzlibrary{positioning, calc, arrows.meta}

\begin{tikzpicture}[
  %% main flow boxes
  mbox/.style={
    rectangle, draw=black, thick, rounded corners=2pt,
    minimum width=6.8cm, text width=6.6cm,
    minimum height=1.3cm, align=center,
    fill=white, font=\small
  },
  %% exclusion boxes (right margin)
  ebox/.style={
    rectangle, draw=black, thick, rounded corners=2pt,
    minimum width=5.0cm, text width=4.8cm,
    minimum height=1.3cm, align=flush left,
    fill=gray!8, font=\small
  },
  %% arrows
  arr/.style={-{Stealth[length=5pt,width=4pt]}, thick},
  %% phase labels
  phase/.style={font=\bfseries\scriptsize, rotate=90, anchor=center}
]

%% ================================================================
%% PHASE 1 — IDENTIFICATION
%% ================================================================
\node[mbox, minimum width=6.0cm, text width=5.8cm] (dbbox) at (-1.8, 0) {
  \textbf{Electronic databases}\\[2pt]
  $n = 5{,}621$\\[2pt]
  {\scriptsize Web of Science: 1,719\\
    Scopus: 1,580\\
    PubMed/MEDLINE: 1,059\\
    EBSCO/CINAHL: 586\\
    Cochrane Library: 399\\
    Google Scholar: 278}
};

\node[mbox, minimum width=5.0cm, text width=4.8cm,
      right=1.5cm of dbbox] (otherbox) {
  \textbf{Other sources (Snowballing)}\\[2pt]
  $n = 7$\\[2pt]
  {\scriptsize (Manual backward reference\\
    tracking)}
};

%% ================================================================
%% PHASE 2 — SCREENING
%% ================================================================
\coordinate (midid) at ($(dbbox.south)!0.5!(otherbox.south)$);

\node[mbox, below=2.6cm of midid, anchor=north] (dedup) {
  \textbf{Database records after deduplication}\\[2pt]
  $n = 3{,}189$\quad{\scriptsize (5,621 $-$ 2,432 duplicates)}
};

%% arrow from databases to deduplication
\draw[arr] (dbbox.south) -- ++(0,-0.8cm) -| (dedup.north);

\node[mbox, below=1.2cm of dedup] (screen1) {
  \textbf{Title screening}\\[2pt]
  $n = 3{,}189$
};

\draw[arr] (dedup) -- (screen1);

\node[ebox, right=2.2cm of screen1] (excl1) {
  \textbf{Excluded by title:} $n = 2{,}680$\\[3pt]
  {\scriptsize
    EC1 (No NLP/ASR of clinical interactions): 1,120\\
    EC2 (No interaction-quality analysis): 584\\
    EC3 (Voice-based diagnosis only): 134\\
    EC4 (Publication type): 281\\
    EC5 (No methodological description): 358\\
    EC6 (Simulated/non-clinical context): 203
  }
};

\draw[arr] (screen1.east) -- (excl1.west);

\node[mbox, below=1.2cm of screen1] (screen2) {
  \textbf{Abstract screening}\\[2pt]
  $n = 509$
};

\draw[arr] (screen1) -- (screen2);

\node[ebox, right=2.2cm of screen2] (excl2) {
  \textbf{Excluded by abstract:} $n = 474$\\[3pt]
  {\scriptsize
    EC1 (No NLP/ASR of clinical interactions): 183\\
    EC2 (No interaction-quality analysis): 116\\
    EC3 (Voice-based diagnosis only): 29\\
    EC4 (Publication type): 43\\
    EC5 (No methodological description): 61\\
    EC6 (Simulated/non-clinical context): 42
  }
};

\draw[arr] (screen2.east) -- (excl2.west);

%% ================================================================
%% PHASE 3 — ELIGIBILITY
%% ================================================================
\node[mbox, below=1.2cm of screen2] (eligib) {
  \textbf{Full text assessed}\\[2pt]
  $n = 35$ (databases) $+$ $7$ (snowballing)
};

\draw[arr] (screen2) -- (eligib);

%% snowballing enters directly at eligibility (bypasses screening)
\draw[arr] (otherbox.south) -- ++(0,-0.8cm) -| ([xshift=1.1cm]eligib.east) -- (eligib.east);

\node[ebox, right=2.2cm of eligib] (excl3) {
  \textbf{Excluded in full text:} $n = 18$ (databases); $0$ (snowballing)\\[3pt]
  {\scriptsize
    Assessed against the full eligibility criteria (EC1--EC6)
  }
};

\draw[arr] (eligib.east) -- (excl3.west);

%% ================================================================
%% PHASE 4 — INCLUDED
%% ================================================================
\node[mbox, below=1.2cm of eligib] (included) {
  \textbf{Studies included in synthesis}\\[2pt]
  $n = 24$\\[2pt]
  {\scriptsize (17 databases + 7 snowballing)}
};

\draw[arr] (eligib) -- (included);

%% ================================================================
%% PHASE LABELS (left margin)
%% ================================================================
\node[phase] at ($(dbbox.west) + (-1.1cm, 0)$)      {IDENTIFICATION};
\node[phase] at ($(dedup.west)  + (-1.1cm, 0)$)      {SCREENING};
\node[phase] at ($(screen1.west) + (-1.1cm, 0)$)      {};   %% screening continues
\node[phase] at ($(screen2.west) + (-1.1cm, 0)$)      {};   %% screening continues
\node[phase] at ($(eligib.west) + (-1.1cm, 0)$)      {ELIGIBILITY};
\node[phase] at ($(included.west) + (-1.1cm, 0)$)    {INCLUDED};

%% ================================================================
%% FOOTNOTE
%% ================================================================
\node[font=\scriptsize, anchor=north west, text width=12cm]
  at ($(included.south west) + (-1.5cm, -0.5cm)$) {%
  Exact counts from the expanded multi-database search (2026).
  Snowballing assessed with the same eligibility criteria.
  $\kappa = 0.78$ (inter-reviewer agreement; Cohen's kappa).
};

\end{tikzpicture}
}
  \caption{Fluxograma PRISMA-ScR do processo de seleção de estudos.}
  \label{fig:prisma}
\end{figure}

% ---------------------------------------------------------------
\subsection{Características dos estudos incluídos}
% ---------------------------------------------------------------

Os 24~estudos do corpus abrangem o período 2014--2025
(Tabela~\ref{tab:estudos}).

\begin{longtable}{p{0.65cm} p{2.4cm} p{1.55cm} p{0.55cm} p{2.2cm} p{2.35cm} p{1.4cm}}
\caption{Características dos estudos incluídos na síntese.}
\label{tab:estudos}\\
\toprule
\textbf{Cód.} & \textbf{Autor, ano} & \textbf{Tipo} &
\textbf{Eixo} & \textbf{Tecnologia} & \textbf{Principal desfecho} & \textbf{País} \\
\midrule
\endfirsthead
\multicolumn{7}{l}{\small\textit{(continuação)}}\\
\toprule
\textbf{Cód.} & \textbf{Autor, ano} & \textbf{Tipo} &
\textbf{Eixo} & \textbf{Tecnologia} & \textbf{Principal desfecho} & \textbf{País} \\
\midrule
\endhead
\bottomrule
\endfoot
%% ---- Eixo A: NLP clínico e extração de entidades biomédicas ----
S01 & Bazoge et al., 2023    & Rev.\ sist.  & A & Transformers, BERT         & NER em CDWs; fenotipagem               & França        \\
S02 & Navarro et al., 2023   & Rev.\ sist.  & A & BERT, BioBERT              & NER clínico (problemas, tratamentos)   & Austrália     \\
S03 & Watabe et al., 2025    & Empírico     & A & BERT-CRF                   & F1\,>\,0{,}8; sintomas relatados       & Japão         \\
S04 & Sezgin et al., 2023    & Empírico     & A & NER + SNOMED CT (zero-shot) & Identificação de sintomas em PGHD     & EUA           \\
S05 & Koleck et al., 2019    & Rev.\ sist.  & A & NLP, classificadores       & Extração de sintomas em EHR            & EUA           \\
S06 & Wang et al., 2018      & Rev.\ lit.   & A & cTAKES, MedLEE, ML         & 263 sistemas de extração clínica       & EUA           \\
%% ---- Eixo B: Documentação clínica automatizada e ASR ----
S07 & Krishna et al., 2021   & Preprint\textsuperscript{‡} & B & CLUSTER2SENT    & Notas SOAP geradas de consultas        & EUA           \\
S08 & Klusty et al., 2025    & Empírico     & B & Whisper + PyAnnote         & Transcrição on-premise; privacidade    & EUA           \\
S09 & Abbasian et al., 2024  & Framework    & B & IA generativa, RAG         & Groundedness, segurança, empatia       & EUA           \\
S10 & Aleksić et al., 2017  & Empírico & B & Resumo de EHR & Dashboard para doença crônica & Sérvia  \\
S11 & Tran et al., 2023      & Empírico     & B & ASR geral vs.\ especializado & WER; erros em pronomes curtos        & EUA           \\
S12 & Schloss \& Konam, 2020 & Preprint\textsuperscript{‡} & B & Bi-LSTM hierárquico & Classificação utterances → SOAP  & EUA           \\
S13 & van Buchem et al., 2021 & Scoping rev. & B & Escribas digitais (revisão) & Lacunas de usabilidade clínica real  & Países Baixos \\
S14 & Falcetta et al., 2023  & Rev.\ sist.  & B & Escribas (revisão)         & Mapeamento de 22 escribas; sem validação prospectiva & Brasil \\
S15 & Ng et al., 2025        & Rev.\ sist.  & B & ASR clínico (revisão)      & WER variou de 8,7\% a $>$50\% conforme contexto & Singapura \\
%% ---- Eixo C: Comunicação médico-paciente ----
S16 & Shao et al., 2025      & Empírico     & C & Escala DPCQ (14 itens)     & Qualidade da comunicação médico-paciente & China       \\
S17 & Muscat et al., 2020    & Framework    & C & Tomada de decisão compartilhada & Literacia em saúde; engajamento & Austrália  \\
S18 & Kelley et al., 2014    & Meta-análise & C & --                         & Relação clínica → desfechos objetivos  & EUA           \\
S19 & Rucinski et al., 2022  & Rev.\ sist.  & C & Métricas de adesão         & Adesão bilateral em ortopedia          & EUA           \\
S20 & McCabe \& Healey, 2018 & Empírico     & C & Análise de conversa        & Autorreparo como preditor de adesão    & RU            \\
S21 & Alkhamees \& Alasqah, 2023 & Rev.\ integr. & C & --               & Barreiras interculturais na comunicação & Arábia S.   \\
%% ---- Eixo D: Biomarcadores vocais e análise multimodal ----
S22 & Fagherazzi et al., 2021 & Framework   & D & Pipeline vocal (acústica+NLP) & Biomarcadores vocais de saúde e emoção & Luxemburgo  \\
S23 & Albert et al., 2024    & Scoping rev.\textsuperscript{‡} & D & IA, Speech, NLP & Comunicação clínico-paciente via IA & França  \\
S24 & Imel et al., 2022      & Relatório    & D & Redes neurais sequenciais  & Valência emocional; concordância humana & EUA          \\
\midrule
\multicolumn{7}{l}{\footnotesize\parbox{13.5cm}{%
  \textsuperscript{‡}Preprint/não peer-reviewed (arXiv/medRxiv) --- incluídos por atenderem IC1--IC4; aceitação de preprints conforme política Elsevier/IJMI.}} \\
\end{longtable}

% ---------------------------------------------------------------
\subsection{Eixo A —NLP clínico e extração de entidades biomédicas}
\label{subsec:eixoA}
% Responsável: Murilo Nogueira Minutti
% ---------------------------------------------------------------

Os estudos do Eixo~A mapeiam a evolução das abordagens de PLN
aplicadas à extração de informação clínica, documentando a transição
de sistemas baseados em regras para arquiteturas de aprendizado
profundo e modelos \textit{transformer}.

O panorama histórico do campo é descrito por \citet{Wang2018}, cuja
revisão de 263~estudos evidencia a predominância de sistemas
simbólicos como cTAKES e MedLEE no processamento de registros
eletrônicos de saúde. Os autores demonstram como a adesão a
abordagens baseadas em regras linguísticas constituiu a norma até a
chegada das técnicas de aprendizado de máquina. A consolidação da
transição para modelos de aprendizado profundo é documentada por
\citet{Bazoge2023}, em revisão sistemática sobre o uso de NLP em
armazéns de dados clínicos (\textit{clinical data warehouses} --
CDWs): os autores evidenciam a migração de métodos simbólicos para
\textit{transformers}, identificando sua aplicabilidade em tarefas de
fenotipagem e gestão de sintomas a partir de dados secundários.

No domínio específico do Reconhecimento de Entidades Nomeadas (NER)
biomédico, \citet{Navarro2023} revisaram sistematicamente a literatura
sobre NER e extração de relações em texto livre médico e concluíram
que modelos como BERT lideram a extração das categorias ``problemas''
e ``tratamentos''. Os autores apontam, contudo, a escassez de
implementações em ambiente clínico real como limitação persistente
do campo.

A capacidade dos modelos \textit{transformer} em capturar a linguagem
coloquial do paciente foi demonstrada por \citet{Watabe2025}, que
desenvolveram e validaram um modelo BERT-CRF para extração de
sintomas relatados em linguagem cotidiana, alcançando F1-score
superior a~0,8. O estudo evidencia a necessidade de ajuste fino
(\textit{fine-tuning}) com dados narrativos autênticos para capturar
gírias e sutilezas linguísticas presentes na fala do paciente.
Complementarmente, \citet{Sezgin2023} demonstraram a viabilidade de
uma abordagem \textit{zero-shot} combinando NER e a ontologia SNOMED~CT
para identificar sintomas em dados gerados por pacientes
(\textit{patient-generated health data} -- PGHD), sem necessidade de
re-treinamento específico, validando a possibilidade de modelos
pré-treinados traduzirem linguagem leiga em métricas clínicas
estruturadas.

Na intersecção entre NLP e prontuário eletrônico, \citet{Koleck2019}
sintetizaram a eficácia de algoritmos de classificação para extração
de sintomas subjetivos -- como dor e alterações de humor -- a partir
de narrativas em texto livre de EHR, evidenciando a viabilidade
técnica da extração automatizada mesmo para dimensões de natureza
subjetiva.

Em conjunto, os estudos do Eixo~A confirmam a maturidade técnica do
NLP clínico para extração de entidades biomédicas, com BERT e suas
variantes estabelecidos como estado da arte. A lacuna transversal
identificada -- apontada por \citet{Navarro2023}
-- é a distância entre desempenho em ambiente laboratorial e a
implantação em sistemas clínicos reais, onde privacidade de dados,
variabilidade linguística e integração com fluxos de trabalho
assistenciais permanecem desafios não resolvidos.

% ---------------------------------------------------------------
\subsection{Eixo B —Documentação clínica automatizada e ASR}
\label{subsec:eixoB}
% Responsável: Fabrício Evangelista dos Santos
% ---------------------------------------------------------------

Os estudos do Eixo~B cobrem o ciclo completo da documentação clínica
automatizada: do reconhecimento da fala à geração de registros
estruturados, incluindo os desafios técnicos e os requisitos éticos
para implantação hospitalar.

As bases da documentação automatizada são estabelecidas por
\citet{Aleksic2017}, cujo modelo de sumarização de prontuários
eletrônicos demonstra a viabilidade de \textit{dashboards} capazes
de identificar pacientes que necessitam de atenção especializada,
validando a síntese legível como estratégia para reduzir a carga
cognitiva do profissional.

No domínio específico da geração de notas SOAP, dois estudos abordam
estratégias complementares. \citet{Schloss2020} utilizaram redes
Bi-LSTM hierárquicas para classificar enunciados de consultas nas
seções SOAP com alta precisão, propondo estratégias modulares para
tratar a ``poluição'' de dados em transcrições automáticas ruidosas.
\citet{Krishna2021}, por sua vez, propuseram o modelo CLUSTER2SENT,
baseado em sumarização modular, para gerar notas estruturadas a partir
de conversas densas entre médico e paciente, facilitando a localização
de evidências verbais pelo profissional na saída analítica.

No que diz respeito à transcrição automática, \citet{Klusty2025}
descreveram um sistema baseado em OpenAI Whisper e PyAnnote operando
em servidores fechados (\textit{on-premise}), resolvendo
simultaneamente o desafio da diarização de locutores e os requisitos
de privacidade para implantação hospitalar. Complementarmente,
\citet{Tran2023} demonstraram que falhas em pronomes curtos
(como \enquote{eu} e \enquote{você}) são mais críticas para a
validade dos indicadores clínicos do que erros em termos médicos
complexos, alertando para a necessidade de precisão na linguagem
casual para garantir a qualidade dos dados extraídos.

As abordagens baseadas em IA generativa trazem um conjunto distinto
de exigências métricas. \citet{Abbasian2024} propõem métricas de
segurança, validade factual (\textit{groundedness}) e empatia como
indicadores obrigatórios para avaliar conversações em saúde geradas
por modelos de linguagem, fundamentando a introdução de Geração
Aumentada por Recuperação (RAG) como mecanismo para ancorar as
respostas do modelo em fatos extraídos da consulta e mitigar
alucinações algorítmicas.

O estado do campo é sintetizado por \citet{vanBuchem2021}, em revisão
de escopo sobre a implantação de escribas digitais na prática clínica:
os autores apontam a ausência de estudos prospectivos de usabilidade
e utilidade clínica real como a lacuna central da área, com a maioria
das iniciativas ainda restrita a provas de conceito sem validação clínica prospectiva.

Em conjunto, os estudos do Eixo~B revelam um campo em acelerada
evolução técnica (com \textit{pipelines} de ASR, diarização e
geração de notas SOAP progressivamente mais robustos), mas com
validação clínica prospectiva ainda escassa \citep{vanBuchem2021}.
A maioria das abordagens direciona o processamento para a automação
documental, subordinando a análise da qualidade da interação clínica
à produção do registro estruturado.

% ---------------------------------------------------------------
\subsection{Eixo C —Comunicação médico-paciente}
\label{subsec:eixoC}
% Responsável: Lucas Kenzo Terazzin Ida
% ---------------------------------------------------------------

Os estudos do Eixo~C estabelecem a base empírica para a qualidade
da comunicação profissional-paciente como determinante de desfechos
clínicos, identificando as dimensões mensuráveis da interação
terapêutica e as barreiras que comprometem sua efetividade.

O impacto da relação clínico-paciente sobre desfechos objetivos de
saúde é demonstrado pela meta-análise de \citet{Kelley2014}, que
comprova o efeito significativo das dimensões emocional e cognitiva
da comunicação sobre resultados clínicos mensuráveis, fornecendo a
fundamentação teórica para o investimento em métricas automáticas de
qualidade da relação terapêutica.

No plano da mensuração estruturada, \citet{Shao2025} desenvolveram e
validaram a escala DPCQ (\textit{Doctor-Patient Communication
Quality}) de 14~itens, com foco em suporte emocional e educação em
saúde. As categorias semânticas propostas pelos autores (interação
emocional versus orientação comportamental) são diretamente
aplicáveis como alvos para métricas automáticas derivadas da fala.

A dimensão da adesão terapêutica em contexto ortopédico é abordada
por \citet{Rucinski2022}, cuja revisão sistemática demonstra que as
métricas vigentes falham ao ignorar as barreiras impostas pela equipe
de saúde e pelo sistema, e não apenas pelo paciente. Os autores
sustentam a criação de dados auditáveis capazes de mapear falhas
bilaterais na reabilitação ortopédica, argumento que reforça a
necessidade de instrumentos de análise da interação, e não apenas
do comportamento isolado do paciente.

Em nível microscópico da fala, \citet{McCabe2018} comprovaram que o
uso de \enquote{autorreparo} por médicos (isto é, a correção
espontânea de enunciados durante a consulta) prediz adesão superior
ao tratamento. O achado valida a análise de dinâmicas linguísticas de
baixo nível como preditores de resultados clínicos, demonstrando que
marcadores comunicacionais identificáveis automaticamente podem
carregar valor prognóstico.

O papel da literacia em saúde na comunicação efetiva é analisado por
\citet{Muscat2020}, que demonstram que a tomada de decisão
compartilhada exige a construção ativa de habilidades no paciente por
meio de instrução verbal eficaz, em contraste com o modelo
informacional passivo. O monitoramento da fala no contexto clínico
deve, portanto, capturar indicadores do sucesso comunicacional da
literacia, e não apenas a transferência unilateral de informação.

Contextos de diversidade linguística e cultural introduzem uma camada
adicional de complexidade. \citet{Alkhamees2023} mapearam como
disparidades linguísticas e estilos autoritários de condução da
consulta geram inibição do paciente e comprometem a qualidade das
decisões clínicas. No contexto diversificado do SUS, essas barreiras
permanecem frequentemente invisíveis aos gestores, mas
potencialmente identificáveis por análise automatizada da interação.

A propedêutica clínica sistematizada por \citet{Porto2019} organiza
a anamnese como sequência estruturada de elementos de inquirição,
definindo os caracteres propedêuticos da queixa: localização,
intensidade, qualidade, irradiação, fatores desencadeantes e de
alívio, duração e evolução, como alvos semânticos obrigatórios do
raciocínio diagnóstico. Essa estrutura constitui a gramática clínica
da interação verbal: o conjunto de informações que o profissional
deve elicitar e o paciente deve verbalizar para que a anamnese
cumpra sua função diagnóstica. Nenhum dos estudos do corpus desta
revisão cruza o modelo propedêutico com a análise automática da
interação verbal: os estudos do Eixo~C mensuram a qualidade da
comunicação em suas dimensões relacionais e cognitivas, mas não
verificam se os caracteres propedêuticos foram efetivamente
verbalizados, compreendidos e registrados durante a consulta.

Em conjunto, os estudos do Eixo~C fornecem o mapa do que deve ser
medido na interação clínica: dimensões emocionais e cognitivas
\citep{Kelley2014,Shao2025}, marcadores linguísticos de adesão
\citep{McCabe2018,Rucinski2022}, competências de literacia
\citep{Muscat2020} e barreiras interculturais \citep{Alkhamees2023}.
A lacuna que emerge é dupla: nenhum estudo operacionaliza essas
dimensões como variáveis derivadas automaticamente da fala
(lacuna tecnológica), e nenhum verifica se os caracteres
propedêuticos da anamnese \citep{Porto2019} foram elicitados com
qualidade comunicacional adequada durante a interação real
(lacuna clínico-metodológica). A mensuração permanece dependente de
escalas de autorrelato ou observação direta, sem integração com
\textit{pipelines} de análise acústico-semântica orientados pela
estrutura clínica da anamnese.

% ---------------------------------------------------------------
\subsection{Eixo D —Biomarcadores vocais e análise multimodal}
\label{subsec:eixoD}
% Responsável: Fabrício Evangelista dos Santos
% ---------------------------------------------------------------

Os estudos do Eixo~D exploram a dimensão acústica e prosódica da
comunicação clínica, estendendo a análise além do conteúdo semântico
do texto transcrito para os sinais vocais que acompanham a fala.

\citet{Fagherazzi2021} descrevem o \textit{pipeline} para captura de
características linguísticas e acústicas, prosódia, qualidade vocal
e ritmo, que indicam estados de saúde e estados emocionais do
locutor. Os autores demonstram que sinais derivados diretamente da
voz podem funcionar como biomarcadores clínicos, fundamentando a
ideia de que a fala do paciente e do profissional, além de seu
conteúdo semântico, carrega informação analítica utilizável
independentemente da transcrição textual.

No domínio da análise de consultas clínicas reais, \citet{Imel2022}
mapearam valências emocionais a partir de transcrições textuais de
diálogos médico-paciente por meio de redes neurais sequenciais,
alcançando concordância com juízes humanos treinados. Embora o estudo
demonstre que marcadores de valência emocional da interação verbal
possam ser identificados automaticamente em nível de enunciado com
confiabilidade comparável à avaliação humana (correlação de 0{,}60),
os autores ressaltam que o modelo foi treinado e avaliado
exclusivamente com base no conteúdo textual dos diálogos (diários de
transcrição), sem processar diretamente o sinal de áudio ou a camada
acústica.

A fronteira não resolvida do Eixo~D é identificada por
\citet{Albert2024}, em revisão de escopo sobre métodos de IA, fala e
NLP para avaliação da comunicação clínico-paciente: os autores
ressaltam a carência de estudos sobre linguagem não verbal, tons de
voz e silêncios. A maioria das abordagens existentes foca em
transcrição e análise semântica, negligenciando a camada
acústico-prosódica como objeto de análise independente. Silêncios,
pausas e modulações de voz, componentes terapêuticos e relacionais
da anamnese, permanecem sem método consolidado de extração e
interpretação automática.

Em conjunto, os estudos do Eixo~D revelam um campo em estado inicial
de desenvolvimento: os fundamentos técnicos de processamento de voz
existem \citep{Fagherazzi2021} e há viabilidade no mapeamento de
marcadores emocionais baseado em transcritos textuais de interações
reais \citep{Imel2022}, mas a integração multimodal que de fato una
os aspectos acústicos da voz aos semânticos do texto na comunicação
clínica permanece incipiente \citep{Albert2024}. Essa lacuna adquire contorno clínico preciso
quando confrontada com a propedêutica de \citet{Porto2019}: se os
caracteres propedêuticos da queixa têm correlatos acústico-prosódicos
na fala (tensão vocal para dor, hesitação e pausas para dúvida,
desengajamento prosódico para barreiras), nenhum estudo do corpus
orienta a extração de biomarcadores vocais pela gramática clínica da
anamnese. A integração entre o modelo propedêutico e o
\textit{pipeline} acústico-semântico representa a fronteira
metodológica mais inexplorada desta revisão.

%% ================================================================
\section{Discussão}
\label{sec:discussao}


\subsection*{Dois eixos paralelos sem convergência sistemática}

Os estudos mapeados revelam dois eixos de desenvolvimento que avançam em
paralelo sem convergência sistemática na literatura.
O primeiro eixo, formado pelos estudos de NLP clínico (Eixo~A) e de
documentação automatizada (Eixo~B), trata a fala como insumo para
documentação: transcreve-se o encontro clínico para gerar notas SOAP,
identificar entidades biomédicas ou popular campos do prontuário eletrônico
\citep{Bazoge2023,vanBuchem2021,Falcetta2023,Ng2025}.
O segundo eixo, formado pelos estudos de comunicação médico-paciente
(Eixo~C) e de biomarcadores vocais (Eixo~D), estabelece o que a comunicação
significa clinicamente e demonstra que sinais acústicos carregam informação
diagnóstica \citep{Kelley2014,Fagherazzi2021,Albert2024}.
Nessa tradição, a qualidade da comunicação é operacionalizada por
instrumentos de autorrelato ou observação presencial, e os biomarcadores
vocais são aplicados majoritariamente à detecção de estados de saúde
(doença de Parkinson, depressão, COVID-19), não à dinâmica comunicacional
durante a consulta.

A heterogeneidade metodológica reflete essa divisão estrutural.
Estudos de NLP como \citet{Bazoge2023} e
\citet{Wang2018} empregam pipelines de extração de
entidades sobre texto retroativo de prontuários; estudos de ASR como
\citet{vanBuchem2021} e
\citet{Klusty2025} medem a acurácia de transcrição
e a redução do tempo de documentação; estudos de comunicação como
\citet{Shao2025} e \citet{Kelley2014}
derivam indicadores relacionais por questionários validados.
Essas tradições compartilham o locus da interação clínica, mas não
compartilham objetos de mensuração, tornando a síntese quantitativa
inviável com os estudos do corpus atual.

\subsection*{Lacuna confirmada: a anamnese como objeto de análise}

Nenhum estudo do corpus operacionaliza a anamnese como evento clínico
estruturado para extração automática de marcadores comunicacionais e
comportamentais derivados da fala.
A distinção é crítica: a maioria dos estudos trata a interação clínica como
\emph{meio}, para produzir documentação ou detectar patologias, não como
\emph{objeto} de análise comunicacional.
\citet{Krishna2021} e
\citet{Schloss2020} classificam enunciados pela
estrutura SOAP, não pela qualidade relacional do encontro.
\citet{Abbasian2024} avaliam a eficácia de
conversas geradas por IA generativa, não o comportamento comunicativo do
profissional de saúde em contexto real.
\citet{Shao2025} e \citet{McCabe2018} descrevem dimensões relevantes de qualidade
comunicativa, clareza, empatia, barreiras à compreensão, mas as
operacionalizam por instrumentos de autorrelato, não por marcadores
automaticamente extraídos da fala.

A revisão de \citet{Albert2024} é a que mais
se aproxima da intersecção: mapeia métodos de SA e NLP para avaliação
da comunicação clínica, mas confirma a escassez de estudos que vinculem
marcadores acústico-prosódicos a desfechos comunicacionais mensuráveis,
especialmente em contextos de reabilitação.
\citet{Porto2019} oferece um referencial semiológico
preciso para a anamnese ortopédica, incluindo hierarquia de sintomas,
cronologia e exploração funcional, que poderia orientar a definição
operacional desses marcadores; no entanto, nenhum estudo recuperado nesta
revisão integra esse referencial com análise automatizada de fala.

\subsection*{Implicações para estudos prospectivos}

A lacuna identificada aponta para um conjunto de requisitos para estudos
futuros: (i)~gravação de encontros clínicos reais em condições ecológicas;
(ii)~ASR com modelos ajustados ao Português Brasileiro;
(iii)~extração de marcadores comunicacionais (expressão de dúvida,
dor, engajamento, barreiras linguísticas) ancorada no referencial
semiológico da anamnese \citep{Porto2019}; e
(iv)~correlação desses marcadores com desfechos de adesão
e qualidade do cuidado em seguimento longitudinal.
\citet{BaenaNeto2026} propõem
exatamente essa integração: um estudo observacional prospectivo em
consultas de reabilitação ortopédica no SUS, com transcrição via ASR
(Whisper \citep{Radford2023}) em processamento \emph{on-premise} e
extração de marcadores para análise longitudinal de desfechos.
Essa abordagem ancora-se nos repertórios técnicos dos Eixos~A, B e D
e nos construtos comunicacionais do Eixo~C, integrando-os de forma inédita
em torno da anamnese como unidade de análise primária.

\subsection*{Limitações desta revisão}

Esta revisão apresenta três limitações centrais.
Primeiro, embora a busca bibliográfica tenha sido amplamente expandida para cobrir múltiplas bases de dados eletrônicas indexadas (PubMed/MEDLINE, Scopus, Web of Science, CINAHL/EBSCO, Cochrane Library e Google Scholar) sem filtros de texto completo gratuito, a inclusão final manteve o corpus de 17 estudos de bases devido à aplicação estrita dos critérios de elegibilidade durante as fases de triagem e leitura na íntegra.
Segundo, há viés de idioma: a literatura mapeada é predominantemente em inglês, com sub-representação de estudos conduzidos no contexto de saúde brasileiro e em Português, o que pode subestimar iniciativas locais de Speech Analytics aplicadas ao SUS.
Terceiro, a heterogeneidade de tipo de estudo, que inclui revisões sistemáticas, estudos empíricos, preprints e modelos conceituais, impede a realização de uma síntese quantitativa ou meta-análise estruturada, limitando também as comparações diretas de qualidade metodológica entre as abordagens.

\subsection*{Considerações éticas}

A aplicação de Speech Analytics em consultas clínicas levanta questões
éticas que extrapolam as considerações padrão para pesquisa com seres
humanos.
A voz constitui dado biométrico sensível nos termos da Lei Geral de
Proteção de Dados Pessoais (LGPD, Lei nº~13.709/2018), impondo obrigações
de minimização da coleta, pseudonimização, separação entre base identificada
e base analítica, e controle de acesso restrito.
A pesquisa deve ser conduzida em conformidade com a Resolução
CNS nº~466/2012 \citep{CNS466_2012} e a Resolução CNS nº~510/2016
\citep{CNS510_2016}, e com a Lei nº~14.874/2024 \citep{Lei14874_2024},
que atualiza o marco regulatório da pesquisa clínica no Brasil.
O Termo de Consentimento Livre e Esclarecido (TCLE) é obrigatório
para a gravação de voz e não admite dispensa quando a fala constitui o
próprio objeto de análise.

Do ponto de vista técnico, \citet{Klusty2025}
demonstram que o processamento \emph{on-premise} é viável em ambiente
hospitalar e representa a solução mais adequada à privacidade, eliminando
o envio de dados sensíveis a servidores externos.
Essa abordagem alinha-se aos cinco princípios para IA responsável
propostos por \citet{Floridi2018} (beneficência,
não-maleficência, autonomia, justiça e explicabilidade), e às
recomendações internacionais da OCDE \citep{OECD2024} e da
Organização Mundial da Saúde \citep{WHO2024} para governança ética de
modelos de IA de grande escala em saúde.
O protocolo de pesquisa deve, ademais, explicitar a proibição de uso
dos marcadores extraídos para decisão clínica autônoma, preservando o
papel decisório do profissional de saúde e a autonomia informada do
paciente.

%% ================================================================
\section{Conclusão}
\label{sec:conclusao}

Esta scoping review mapeou o estado da arte em Speech Analytics e
comunicação clínica com o objetivo de estabelecer a fundamentação
científica para estudos prospectivos que integrem análise automatizada
da fala à avaliação da qualidade comunicacional em contextos de
reabilitação ortopédica.

Os quatro eixos temáticos evidenciam avanços consolidados em NLP clínico
e ASR: pipelines de extração de entidades biomédicas alcançam desempenho
competitivo em prontuários eletrônicos \citep{Bazoge2023,Koleck2019},
modelos de reconhecimento de fala como o Whisper \citep{Radford2023}
viabilizam a transcrição de encontros clínicos em ambiente hospitalar
\citep{Klusty2025,Ng2025}, e biomarcadores vocais demonstram capacidade
de capturar estados fisiológicos e emocionais relevantes para o cuidado
em saúde \citep{Fagherazzi2021,Albert2024}.
A literatura de comunicação médico-paciente, por sua vez, estabelece
com consistência que a qualidade relacional do encontro clínico influencia
a adesão terapêutica e os desfechos de saúde \citep{Kelley2014,Rucinski2022}.

No entanto, a revisão confirma uma lacuna consistente e específica:
poucos estudos operacionalizam interações clínicas como fonte de
variáveis mensuráveis derivadas da fala, especialmente marcadores
comunicacionais e comportamentais (expressão de dúvida, dor, engajamento,
barreiras à compreensão) extraídos automaticamente durante a anamnese
como evento estruturado.
A maioria das iniciativas trata a fala como insumo para documentação ou
como sinal para detecção de patologias, sem ancorar a análise no
referencial semiológico da anamnese \citep{Porto2019}, que define
o que, na fala, é clinicamente significativo do ponto de vista
comunicacional.

Esta lacuna aponta para a necessidade de estudos prospectivos que integrem
Speech Analytics à análise estruturada da comunicação clínica para
monitoramento de adesão e qualidade do cuidado em tempo real.
Tais estudos requerem, além de infraestrutura técnica adequada, um
protocolo ético robusto que trate a voz como dado biométrico sensível,
garanta o consentimento informado e preserve a privacidade dos pacientes
por meio de processamento local dos dados.
O protocolo metodológico descrito em \citet{BaenaNeto2026} representa uma resposta direta a esse
conjunto de requisitos no contexto do Sistema Único de Saúde.

%% ================================================================
\section*{Declarações}

\subsection*{Financiamento}
Este estudo não recebeu financiamento de agências de fomento ou entidades
dos setores público, comercial ou sem fins lucrativos.
Os custos de processamento do artigo (APC) são dispensados para o autor
correspondente, afiliado à Universidade São Francisco (USF, Bragança
Paulista, SP, Brasil), ao abrigo do Acordo de Acesso Aberto firmado entre
a Coordenação de Aperfeiçoamento de Pessoal de Nível Superior (CAPES)
e a Elsevier para o período 2026--2028
(\textit{Read \& Publish Agreement}).

\subsection*{Conflito de interesses}
Nenhum dos autores possui vínculo de emprego, consultoria, participação
acionária, honorários, patentes ou financiamento por entidades com
interesse comercial ou institucional no tema desta revisão.
O protocolo de estudo prospectivo citado no texto
\citep{BaenaNeto2026} é de autoria parcial do autor correspondente
(R.B.N.) e do coautor sênior (V.I.B.S.); essa relação é integralmente
divulgada pela lista de autoria e não constitui conflito de interesses
nos termos das diretrizes do ICMJE\@.

\subsection*{Aprovação ética}
Não aplicável. Este estudo consiste em uma revisão bibliográfica baseada
exclusivamente em literatura científica publicada e acessível ao público.
Não foram coletados dados primários, não foram recrutados participantes e
nenhuma intervenção foi realizada sobre seres humanos.
Nos termos da Resolução CNS nº~510/2016 \citep{CNS510_2016}, pesquisas
que utilizam informações de domínio público estão dispensadas de registro e
avaliação pelo sistema CEP/CONEP\@.

\subsection*{Disponibilidade de dados}
Os dados que suportam os achados desta revisão estão disponíveis da
seguinte forma:
(i)~as \textbf{strings completas de busca} por base de dados estão
registradas no arquivo suplementar~S1 e no repositório público do projeto
(\url{https://doi.org/10.5281/zenodo.20821063});
(ii)~o \textbf{formulário de extração de dados} padronizado (campos
extraídos por estudo, incluindo autores, ano, delineamento, população,
tecnologia utilizada, desfechos e avaliação de qualidade metodológica)
está disponível como material suplementar~S2;
(iii)~as \textbf{decisões de triagem} (inclusão/exclusão por título,
resumo e texto completo), com os respectivos motivos de exclusão,
foram registradas na plataforma Rayyan \citep{Ouzzani2016} e são
disponibilizadas mediante solicitação ao autor correspondente;
(iv)~a \textbf{lista completa de estudos incluídos}, com classificação
por eixo temático e pontuação de qualidade, integra a Tabela~\ref{tab:estudos}
no corpo do artigo.
Os dados brutos de gravações, transcrições ou registros de pacientes
não se aplicam a este estudo, que se baseia exclusivamente em literatura
publicada.

\subsection*{Uso de inteligência artificial}
Em conformidade com a Política de IA Generativa da Elsevier, os autores
declaram que os seguintes modelos de inteligência artificial foram
utilizados na preparação deste trabalho:
(i)~\textbf{Claude Sonnet~4.6} (Anthropic PBC,
\url{https://claude.ai}, 2025--2026), empregado como assistente nas buscas
bibliográficas nas bases de dados eletrônicas;
(ii)~\textbf{Claude Code} (Anthropic PBC,
\url{https://claude.ai/code}, 2026), utilizado para apoio à triagem do
corpus, extração de dados e tradução de trechos de artigos consultados em
idioma estrangeiro.
Todo o conteúdo intelectual, as análises, as sínteses e as conclusões
são de responsabilidade exclusiva dos autores, que revisaram e validaram
integralmente todo o material descrito produzido com auxílio de IA\@.

\printcredits

%% ================================================================
\bibliographystyle{cas-model2-names}
\bibliography{references}

\end{document}